% ============================================================
\section{Introducción}\label{sec:intro}
% ============================================================

El presente trabajo integrador documenta el desarrollo de \tails{} (\emph{Technical
Articulated Intelligent Linkage System}), un sistema mecatrónico compuesto por un brazo
robótico articulado de tres grados de libertad (3-GDL) con un efector final tipo pinza
(\emph{gripper}), una placa de control propia basada en el microcontrolador
\textbf{STM32F103C8T6} (placa de desarrollo \emph{Blue Pill}) y una interfaz gráfica de
escritorio para su operación desde una PC.

\paragraph{Idea y motivación.}
La consigna de la cátedra, planteada por el Prof.\ Iriarte, propuso desarrollar un robot
manipulador con distintos modos operativos: (1)~calibrar los ejes; (2)~guardar posiciones
articulares o trayectorias, aprendidas manualmente o cargadas por un puerto serie; y
(3)~ejecutar esas posiciones o trayectorias. La propuesta incluía un conjunto de comandos de
ejemplo —\cmd{:Att} para entrar al modo aprendizaje de la trayectoria \texttt{tt}, \cmd{:G}
para guardar la posición actual como siguiente punto, \cmd{:Guuu,vvv,www} para guardar una
posición articular arbitraria, \cmd{:Z} para terminar el aprendizaje, \cmd{:Ett} para ejecutar
una trayectoria aprendida, \cmd{:Duuu,vvv,www} para ejecutar un punto directo y \cmd{:vvvv} para
la velocidad global de 10 a 100~\%— y dejaba explícitamente abierta la posibilidad de replantear
la funcionalidad, por ejemplo que cada punto aprendido incluyera una velocidad local, un tiempo
de espera antes del siguiente movimiento y el estado del \emph{gripper}. El foco del trabajo
está puesto en la \textbf{implementación sobre el
microcontrolador}: configuración de periféricos (temporizadores, interrupciones externas,
PWM, USB e I\textsuperscript{2}C), generación de pulsos para los drivers de motores paso a
paso, gestión de la seguridad (finales de carrera y parada de emergencia) y un protocolo de
comandos robusto entre la PC y el robot. Por este motivo, la cinemática del brazo no se
aborda: el sistema trabaja en coordenadas articulares expresadas en pasos de motor.

\paragraph{Punto de partida.}
La estructura mecánica se basa en el diseño abierto \emph{Robot Arm MK2 Plus (Stepper Motor
in Used)} publicado por el usuario \emph{jackyltle} en Thingiverse en septiembre de
2017~\cite{thingiverse}, pensado para imprimirse en 3D por deposición de material fundido
(FDM). El diseño original se controla con un Arduino UNO (ATmega328P), una \emph{shield} CNC
y drivers Pololu A4988 para tres motores NEMA~17, y no contempla sensores de fin de carrera ni
alojamientos para ellos. Sobre esa base se diseñaron soportes propios para los sensores, una
carcasa para la PCB, una placa de control dedicada en KiCad y un firmware completamente nuevo
para la arquitectura ARM Cortex-M3.

\fig[0.75\textwidth]{mk2_original}{Brazo \emph{Robot Arm MK2 Plus} original, tomado como base mecánica~\cite{thingiverse}.}{fig:mk2}{render o foto del diseño original (carpeta \texttt{diseño 3D/jackyltle\_Structure/images})}

\paragraph{Utilidad.}
El sistema permite: (i) referenciar los ejes de forma automática (\emph{homing}) contra
finales de carrera; (ii) mover el brazo manualmente (\emph{jogging}) desde la PC o el teclado
y memorizar puntos con velocidad, espera y estado de pinza; (iii) guardar esas secuencias en
archivos JSON y reproducirlas automáticamente, y (iv) detener el sistema en forma inmediata
ante una emergencia, por botón físico o por software. Es una plataforma didáctica y de
prototipado para tareas simples de \emph{pick and place}.

\paragraph{Principios de funcionamiento.}
Cada articulación es accionada por un motor paso a paso gobernado por un driver A4988 con
interfaz STEP/DIR. El microcontrolador genera los trenes de pulsos a partir de una
interrupción periódica de \SI{20}{\kilo\hertz}, aplicando un perfil de velocidad
trapezoidal. La pinza utiliza un servomotor SG90 comandado por PWM a \SI{50}{\hertz}. La
posición de cada eje se conoce en lazo abierto contando los pasos emitidos, referidos al cero
que fija el \emph{homing}. La PC y el robot se comunican por USB (clase CDC, puerto COM
virtual) mediante tramas ASCII, y el robot devuelve telemetría periódica de su estado.

\fig[0.8\textwidth]{robot_general}{Vista general del brazo \tails{} con su electrónica de control.}{fig:robot-general}{foto del robot completo, p.~ej. \texttt{Documentación/Fotos/extra/TAILS.jpg}}
